% ============================================================
\chapter{It\^o Stochastic Integral}
\label{ch:ito_integral}
% ============================================================

\section{Motivation}

In 1942, in Kyoto, the mathematician Kiyosi It\^o tackled a fundamental problem: how does one define an integral with respect to Brownian motion? Brownian paths are continuous but nowhere differentiable, and their total variation is infinite on every interval. The classical Stieltjes integral simply does not apply. It\^o circumvented the obstacle through an ingenious construction: he approximated the integrand by step processes and showed that the limit exists in $L^2$, thanks to the It\^o isometry. The choice of evaluating at the \emph{left endpoint} of each subinterval (unlike Stratonovich, who uses the midpoint) gives the It\^o integral the property of being a martingale --- a property crucial for finance and probability.

The goal of this chapter is to give rigorous meaning to the integral
\[
  \int_0^t H_s \, dB_s,
\]
where $(B_t)$ is a Brownian motion and $(H_t)$ an adapted process. Since
Brownian sample paths have infinite total variation
(Theorem~\ref{thm:total_var}), this integral cannot be defined in the Stieltjes
sense. It\^o's construction proceeds by approximation.

\section{Construction for Simple Processes}

\begin{definition}[Simple Process]
A process $(H_t)_{t \in [0,T]}$ is \emph{simple} (or \emph{elementary}) if there
exist a partition $0 = t_0 < t_1 < \cdots < t_n = T$ and bounded
$\mathcal{F}_{t_k}$-measurable r.v.\ $H_k$ such that:
\[
  H_t = H_0 \mathbf{1}_{\{0\}}(t) + \sum_{k=0}^{n-1} H_k \mathbf{1}_{(t_k, t_{k+1}]}(t).
\]
We write $\mathcal{S}$ for the space of simple processes.
\end{definition}

\begin{definition}[It\^o Integral for Simple Processes]
For $H \in \mathcal{S}$, we define:
\[
  I(H) = \int_0^T H_s \, dB_s := \sum_{k=0}^{n-1} H_k (B_{t_{k+1}} - B_{t_k}).
\]
\end{definition}

\begin{warning}[title=\textbf{Left-Point Evaluation}]
In It\^o's definition, the integrand $H_k$ is evaluated at the \emph{left}
endpoint $t_k$ of the interval $(t_k, t_{k+1}]$. This choice is crucial for
the martingale property. A different choice (e.g.\ the midpoint, as in the
Stratonovich integral) gives a different result.
\end{warning}

\begin{proposition}[Properties for Simple Processes]
\label{prop:simple_props}
For $H, G \in \mathcal{S}$ and $a, b \in \R$:
\begin{enumerate}
  \item \textbf{Linearity}: $I(aH + bG) = aI(H) + bI(G)$.
  \item \textbf{Zero expectation}: $\E[I(H)] = 0$.
  \item \textbf{It\^o isometry}:
    $\E[I(H)^2] = \E\bigl[\int_0^T H_s^2 \, ds\bigr]$.
  \item \textbf{Martingale property}: the process
    $t \mapsto \int_0^t H_s \, dB_s$ is a martingale.
\end{enumerate}
\end{proposition}

\begin{proof}[Proof of It\^o isometry]
\begin{align*}
  \E[I(H)^2] &= \E\Bigl[\Bigl(\sum_{k=0}^{n-1} H_k \Delta B_k\Bigr)^2\Bigr] \\
  &= \sum_{k=0}^{n-1} \E[H_k^2 (\Delta B_k)^2]
    + 2 \sum_{j < k} \E[H_j \Delta B_j H_k \Delta B_k].
\end{align*}
The cross terms vanish since, for $j < k$:
\[
  \E[H_j \Delta B_j H_k \Delta B_k]
  = \E[\underbrace{H_j \Delta B_j H_k}_{\mathcal{F}_{t_k}\text{-measurable}}
  \underbrace{\E[\Delta B_k \mid \mathcal{F}_{t_k}]}_{=0}] = 0.
\]
For the diagonal terms:
\[
  \E[H_k^2 (\Delta B_k)^2] = \E[H_k^2 \E[(\Delta B_k)^2 \mid \mathcal{F}_{t_k}]]
  = \E[H_k^2 (t_{k+1} - t_k)].
\]
Hence $\E[I(H)^2] = \sum_k \E[H_k^2 (t_{k+1} - t_k)] = \E[\int_0^T H_s^2 \, ds]$.
\end{proof}

\section{Extension to $L^2$ Processes}

\begin{definition}[Space $\mathcal{H}^2$]
The space of admissible integrands is:
\[
  \mathcal{H}^2 = \mathcal{H}^2([0, T]) = \Bigl\{H \text{ adapted}
  : \E\Bigl[\int_0^T H_s^2 \, ds\Bigr] < \infty \Bigr\}.
\]
Equipped with the norm
$\norm{H}_{\mathcal{H}^2} = \bigl(\E[\int_0^T H_s^2 \, ds]\bigr)^{1/2}$,
this is a Hilbert space.
\end{definition}

\begin{theorem}[Extension of the It\^o Integral]
\label{thm:ito_extension}
The linear isometric map $I : \mathcal{S} \to L^2(\Omega)$ extends uniquely to
an isometry $I : \mathcal{H}^2 \to L^2(\Omega)$. The extended integral, denoted
\[
  \int_0^T H_s \, dB_s = I(H),
\]
satisfies the same properties as for simple processes: linearity, zero
expectation, It\^o isometry.
\end{theorem}

\begin{proof}[Proof sketch]
The space $\mathcal{S}$ is dense in $\mathcal{H}^2$ (any adapted
square-integrable process can be approximated by simple processes). The It\^o
isometry guarantees that $I$ is an isometry from
$(\mathcal{S}, \norm{\cdot}_{\mathcal{H}^2})$ into
$(L^2(\Omega), \norm{\cdot}_{L^2})$. By completeness of $L^2(\Omega)$, $I$
extends uniquely.
\end{proof}

\section{Properties of the It\^o Integral}

\begin{theorem}[Fundamental Properties]
\label{thm:ito_props}
Let $H \in \mathcal{H}^2([0,T])$. The process $M_t = \int_0^t H_s \, dB_s$
satisfies:
\begin{enumerate}
  \item \textbf{Continuity}: $(M_t)$ has a version with continuous sample paths.
  \item \textbf{Martingale}: $(M_t)_{0 \leq t \leq T}$ is a square-integrable
    martingale with respect to $(\mathcal{F}_t)$.
  \item \textbf{It\^o isometry}: $\E[M_t^2] = \E[\int_0^t H_s^2 \, ds]$.
  \item \textbf{Quadratic variation}: $\langle M \rangle_t = \int_0^t H_s^2 \, ds$.
\end{enumerate}
\end{theorem}

\begin{keyformulas}[title=\textbf{Key Properties of the It\^o Integral}]
For $M_t = \int_0^t H_s \, dB_s$ with $H \in \mathcal{H}^2$:
\begin{align*}
  \E[M_t] &= 0 \\
  \E[M_t^2] &= \E\Bigl[\int_0^t H_s^2 \, ds\Bigr]
    \quad \text{(It\^o isometry)} \\
  \langle M \rangle_t &= \int_0^t H_s^2 \, ds \\
  dM_t &= H_t \, dB_t
\end{align*}
\end{keyformulas}

\section{Extension to Locally $L^2$ Integrands}

\begin{definition}[Generalised It\^o Integral]
For an adapted process $H$ with $\int_0^T H_s^2 \, ds < \infty$ a.s.\ (but not
necessarily $\E[\int_0^T H_s^2 \, ds] < \infty$), the It\^o integral is defined by
localisation: set $\tau_n = \inf\{t : \int_0^t H_s^2 \, ds \geq n\}$ and
$\int_0^t H_s \, dB_s = \lim_{n \to \infty} \int_0^{t \wedge \tau_n} H_s \, dB_s$.

The resulting process is a continuous \emph{local martingale}.
\end{definition}

\section{Stratonovich Integral}

\begin{definition}[Stratonovich Integral]
The \emph{Stratonovich integral} is defined by:
\[
  \int_0^T H_s \circ dB_s = \lim_{\abs{\Pi} \to 0}
  \sum_{k=0}^{n-1} \frac{H_{t_k} + H_{t_{k+1}}}{2} (B_{t_{k+1}} - B_{t_k}).
\]
\end{definition}

\begin{proposition}[It\^o--Stratonovich Relation]
If $H_t = f(B_t)$ with $f \in C^1(\R)$, then:
\[
  \int_0^T f(B_s) \circ dB_s = \int_0^T f(B_s) \, dB_s
  + \frac{1}{2} \int_0^T f'(B_s) \, ds.
\]
\end{proposition}

\begin{remark}
The Stratonovich integral obeys the classical rules of differential calculus (no
It\^o correction term) but is not a martingale in general. The It\^o integral is
preferred in finance (martingale property), while the Stratonovich integral is
natural in physics (coordinate invariance).
\end{remark}

\section{Computational Examples}

\begin{example}[$\int_0^t B_s \, dB_s$]
For a partition $\Pi = \{0 = t_0 < t_1 < \cdots < t_n = t\}$:
\begin{align*}
  \sum_{k=0}^{n-1} B_{t_k}(B_{t_{k+1}} - B_{t_k})
  &= \sum_{k=0}^{n-1} \Bigl[\frac{1}{2}(B_{t_{k+1}}^2 - B_{t_k}^2)
    - \frac{1}{2}(B_{t_{k+1}} - B_{t_k})^2\Bigr] \\
  &= \frac{1}{2} B_t^2 - \frac{1}{2} \sum_{k=0}^{n-1} (\Delta B_k)^2.
\end{align*}
Taking the limit and using $\sum (\Delta B_k)^2 \to t$ (quadratic variation):
\[
  \int_0^t B_s \, dB_s = \frac{1}{2} B_t^2 - \frac{1}{2} t.
\]
\end{example}

\begin{warning}[title=\textbf{The $-t/2$ Term}]
The result $\int_0^t B_s \, dB_s = \frac{1}{2}B_t^2 - \frac{1}{2}t$ shows that
stochastic calculus differs from ordinary calculus. In ordinary calculus, one
would have $\int_0^t x \, dx = x^2/2$. The term $-t/2$ is the ``It\^o correction
term'', a consequence of the non-zero quadratic variation of Brownian motion.
\end{warning}

\begin{example}[$\int_0^t e^{B_s - s/2} \, dB_s$]
Set $M_t = e^{B_t - t/2}$. One verifies in the next chapter (It\^o's formula)
that $dM_t = M_t \, dB_t$, hence:
\[
  \int_0^t e^{B_s - s/2} \, dB_s = M_t - M_0 = e^{B_t - t/2} - 1.
\]
\end{example}

\section{Exercises}

\begin{exercise}
Compute $\E[(\int_0^1 B_s \, dB_s)^2]$ using the It\^o isometry.
\end{exercise}

\begin{exercise}
Let $f \in C^1(\R)$. Show that $\int_0^t f(s) \, dB_s$ is a Gaussian r.v.\ and
compute its law.
\end{exercise}

\begin{exercise}
Compute $\int_0^t s \, dB_s$ and verify that it is a Gaussian r.v.\ with
variance $t^3/3$.
\end{exercise}

\begin{exercise}
Show that $\int_0^t B_s^2 \, dB_s = \frac{1}{3} B_t^3 - \int_0^t B_s \, ds$
using It\^o's formula (anticipate the next chapter, or prove directly by
approximation).
\end{exercise}

\begin{exercise}
Let $H_t = \text{sgn}(B_t)$ (sign of Brownian motion). Verify that
$H \in \mathcal{H}^2$ and that
$\int_0^t \text{sgn}(B_s) \, dB_s = \abs{B_t} - L_t^0$ (Tanaka's formula).
\end{exercise}

\begin{exercise}
Compare the It\^o and Stratonovich integrals of $\int_0^t B_s \, dB_s$ and
verify the correction formula.
\end{exercise}
